\documentclass[11pt,a4paper]{article}

% --- Packages ---
\usepackage[utf8]{inputenc}
\usepackage[margin=1in]{geometry}
\usepackage{booktabs}
\usepackage{longtable}
\usepackage{listings}
\usepackage{xcolor}
\usepackage{hyperref}
\usepackage{helvet}
\renewcommand{\familydefault}{\sfdefault}

% --- Syntax Highlighting Colors ---
\definecolor{codegreen}{rgb}{0,0.6,0}
\definecolor{codegray}{rgb}{0.5,0.5,0.5}
\definecolor{codepurple}{rgb}{0.58,0,0.82}
\definecolor{backcolour}{rgb}{0.95,0.95,0.92}

% --- Listings Style ---
\lstset{
    backgroundcolor=\color{backcolour},   
    commentstyle=\color{codegreen},
    keywordstyle=\color{magenta},
    numberstyle=\tiny\color{codegray},
    stringstyle=\color{codepurple},
    basicstyle=\ttfamily\footnotesize,
    breakatwhitespace=false,         
    breaklines=true,                 
    captionpos=b,                    
    keepspaces=true,                 
    numbers=left,                    
    numbersep=5pt,                  
    showspaces=false,                
    showstringspaces=false,
    showtabs=false,                  
    tabsize=2,
    frame=single
}

% --- Document Metadata ---
\title{
    \rule{\textwidth}{1pt} \\
    \textbf{\Large Test Data Set} \\
    \large IBA Campus Facility Booking System \\
    \rule{\textwidth}{1pt}
}
\author{
    \textbf{Team Members:} \\
    Abdul Wasay Imran (27126), Saad Imam (27079), \\
    Muhammad Imad Raza (26953), Syed Muhammad Deebaj Kazmi (26012)
}
\date{Version 1.1 --- \today}

\begin{document}

\maketitle
\tableofcontents
\newpage

\section{Overview}
This document defines the data environment used for the Quality Assurance lifecycle of the IBA Campus Facility Booking System. To ensure test repeatability and prevent "flaky" tests, the system utilizes a State Reset Strategy. 

The environment is completely wiped and re-seeded before every E2E test suite and during specific integration tests. This ensures that every test starts from a known, clean state.

\section{Environment Configuration}
The testing environment is isolated from production using a dedicated Supabase PostgreSQL instance.
\begin{itemize}
    \item \textbf{Database Engine:} PostgreSQL (via Supabase)
    \item \textbf{Reset Method:} Drop and Recreate Schema (\texttt{init.sql})
    \item \textbf{Seeding Method:} Static UUID Injection (\texttt{seedTestData.sql})
    \item \textbf{Automation Helper:} TypeScript / Node.js \texttt{pg} client
\end{itemize}

\section{Test Data Summary}
The following tables represent the static data injected into the system during the seeding process. These IDs are hardcoded in our SQL scripts to allow E2E tests (Playwright/Selenium) to target specific entities consistently.

\subsection{Test User Accounts}
\textit{Note: All test accounts use the password: \textbf{testpass}}
\begin{center}
\begin{tabular}{@{}llll@{}}
\toprule
\textbf{ERP / Username} & \textbf{Name} & \textbf{Role} & \textbf{Email} \\ \midrule
test-student & Test Student & Student & student@test.com \\
test-po & Test PO & Program Office & po@test.com \\
test-admin & Test Admin & Admin & admin@test.com \\
test-student-2 & Concurrency Student & Student & student2@test.com \\ \bottomrule
\end{tabular}
\end{center}

\subsection{Facility Data (Buildings \& Rooms)}
Static UUIDs are used to ensure API calls in integration tests remain valid across environment resets.
\begin{center}
\begin{longtable}{@{}lll@{}}
\toprule
\textbf{Entity Name} & \textbf{Type} & \textbf{Static UUID Reference} \\ \midrule
Test Building & Building & \texttt{99999999-9999-4999-8999-999999999999} \\
Test Room A & Room & \texttt{aaaaaaaa-aaaa-4aaa-8aaa-aaaaaaaaaaaa} \\
Test Room B & Room & \texttt{bbbbbbbb-bbbb-4bbb-8bbb-bbbbbbbbbbbb} \\ \bottomrule
\end{longtable}
\end{center}

\section{Seeding \& Reset Logic}
The system uses a two-step process to maintain data integrity during testing:
\begin{enumerate}
    \item \textbf{Global Setup:} The \texttt{global-setup.ts} script triggers a database reset before the test runner (Playwright) starts.
    \item \textbf{Atomic Reset:} The \texttt{resetDatabase()} function executes the SQL scripts via a PostgreSQL client, ensuring the \texttt{bookings} table is cleared and the \texttt{time\_slots} are restored.
\end{enumerate}

\newpage
\appendix
\section{Appendix: Source Scripts}

\subsection{init.sql (Full Schema Setup)}
This script is used to initialize the database from scratch, creating all tables, enums, and static time slots.
\begin{lstlisting}[language=SQL]
-- Full Reset: Drop everything first
DROP TABLE IF EXISTS bookings cascade;


DROP TABLE IF EXISTS blocked_slots cascade;


DROP TABLE IF EXISTS rooms cascade;


DROP TABLE IF EXISTS buildings cascade;


DROP TABLE IF EXISTS users cascade;


DROP TABLE IF EXISTS time_slots cascade;


DROP TYPE if EXISTS user_role cascade;


DROP TYPE if EXISTS booking_status cascade;


DROP TYPE if EXISTS room_type cascade;


-- ============================================================
--  IBA FACILITY BOOKING SYSTEM — Supabase SQL Schema
--  Run this entire file in: Supabase → SQL Editor → New Query
-- ============================================================
-- ── ENUMS ────────────────────────────────────────────────────
CREATE TYPE user_role AS ENUM('student', 'programoffice', 'admin');


CREATE TYPE booking_status AS ENUM('pending', 'approved', 'rejected', 'cancelled');


CREATE TYPE room_type AS ENUM(
    'Classroom',
    'Seminar Hall',
    'Computer Lab',
    'Meeting Room'
);


-- ── USERS ────────────────────────────────────────────────────
CREATE TABLE users (
    id UUID PRIMARY KEY DEFAULT GEN_RANDOM_UUID(),
    erp VARCHAR(50) UNIQUE NOT NULL, -- ERP ID for students, username for others
    name VARCHAR(150) NOT NULL,
    email VARCHAR(150) UNIQUE NOT NULL,
    password VARCHAR(255) NOT NULL, -- bcrypt hashed
    role user_role NOT NULL DEFAULT 'student',
    created_at TIMESTAMPTZ NOT NULL DEFAULT NOW(),
    updated_at TIMESTAMPTZ NOT NULL DEFAULT NOW()
);


-- ── BUILDINGS ────────────────────────────────────────────────
CREATE TABLE buildings (
    id UUID PRIMARY KEY DEFAULT GEN_RANDOM_UUID(),
    name VARCHAR(150) NOT NULL,
    location VARCHAR(200),
    created_at TIMESTAMPTZ NOT NULL DEFAULT NOW()
);


-- ── ROOMS ────────────────────────────────────────────────────
CREATE TABLE rooms (
    id UUID PRIMARY KEY DEFAULT GEN_RANDOM_UUID(),
    building_id UUID NOT NULL REFERENCES buildings (id) ON DELETE CASCADE,
    name VARCHAR(100) NOT NULL,
    capacity INTEGER NOT NULL CHECK (capacity > 0),
    type room_type NOT NULL DEFAULT 'Classroom',
    created_at TIMESTAMPTZ NOT NULL DEFAULT NOW()
);


-- ── TIME SLOTS (static reference table) ─────────────────────
CREATE TABLE time_slots (
    id INTEGER PRIMARY KEY,
    start_time VARCHAR(5) NOT NULL, -- e.g. "08:30"
    end_time VARCHAR(5) NOT NULL, -- e.g. "09:45"
    label VARCHAR(30) NOT NULL -- e.g. "8:30 - 9:45"
);


INSERT INTO
    time_slots (id, start_time, end_time, label)
VALUES
    (1, '08:30', '09:45', '8:30 - 9:45'),
    (2, '10:00', '11:15', '10:00 - 11:15'),
    (3, '11:30', '12:45', '11:30 - 12:45'),
    (4, '13:00', '14:15', '1:00 - 2:15'),
    (5, '14:30', '15:45', '2:30 - 3:45'),
    (6, '16:00', '17:15', '4:00 - 5:15'),
    (7, '17:30', '18:45', '5:30 - 6:45');


-- ── BOOKINGS ─────────────────────────────────────────────────
CREATE TABLE bookings (
    id UUID PRIMARY KEY DEFAULT GEN_RANDOM_UUID(),
    user_id UUID NOT NULL REFERENCES users (id) ON DELETE CASCADE,
    room_id UUID NOT NULL REFERENCES rooms (id) ON DELETE CASCADE,
    slot_id INTEGER NOT NULL REFERENCES time_slots (id),
    date date NOT NULL,
    purpose TEXT NOT NULL,
    status booking_status NOT NULL DEFAULT 'pending',
    reviewed_by UUID REFERENCES users (id), -- PO or admin who actioned it
    created_at TIMESTAMPTZ NOT NULL DEFAULT NOW(),
    updated_at TIMESTAMPTZ NOT NULL DEFAULT NOW(),
    -- Prevent double-booking same room+date+slot
    UNIQUE (room_id, date, slot_id)
);


-- ── BLOCKED SLOTS ────────────────────────────────────────────
CREATE TABLE blocked_slots (
    id UUID PRIMARY KEY DEFAULT GEN_RANDOM_UUID(),
    room_id UUID NOT NULL REFERENCES rooms (id) ON DELETE CASCADE,
    slot_id INTEGER NOT NULL REFERENCES time_slots (id),
    date date NOT NULL,
    reason VARCHAR(200),
    blocked_by UUID NOT NULL REFERENCES users (id),
    created_at TIMESTAMPTZ NOT NULL DEFAULT NOW(),
    UNIQUE (room_id, date, slot_id)
);


-- ── UPDATED_AT TRIGGER ───────────────────────────────────────
CREATE OR REPLACE FUNCTION set_updated_at () returns trigger AS $$
BEGIN
  NEW.updated_at = NOW();
  RETURN NEW;
END;
$$ language plpgsql;


CREATE TRIGGER trg_users_updated_at before
UPDATE ON users FOR each ROW
EXECUTE function set_updated_at ();


CREATE TRIGGER trg_bookings_updated_at before
UPDATE ON bookings FOR each ROW
EXECUTE function set_updated_at ();


-- ── INDEXES ──────────────────────────────────────────────────
CREATE INDEX idx_bookings_user_id ON bookings (user_id);


CREATE INDEX idx_bookings_room_id ON bookings (room_id);


CREATE INDEX idx_bookings_date ON bookings (date);


CREATE INDEX idx_bookings_status ON bookings (status);


CREATE INDEX idx_blocked_room_date ON blocked_slots (room_id, date);


CREATE INDEX idx_rooms_building ON rooms (building_id);


-- ── SEED DATA ────────────────────────────────────────────────
-- Admin user (password: admin123)
INSERT INTO
    users (erp, name, email, password, role)
VALUES
    (
        'admin',
        'System Admin',
        'admin@iba.edu.pk',
        '$2a$10$92IXUNpkjO0rOQ5byMi.Ye4oKoEa3Ro9llC/.og/at2.uheWG/igi',
        'admin'
    );


-- Program Office (password: password)
INSERT INTO
    users (erp, name, email, password, role)
VALUES
    (
        'po001',
        'Dr. Fatima Zahra',
        'fatima@iba.edu.pk',
        '$2a$10$92IXUNpkjO0rOQ5byMi.Ye4oKoEa3Ro9llC/.og/at2.uheWG/igi',
        'programoffice'
    );


-- Students (password: password)
INSERT INTO
    users (erp, name, email, password, role)
VALUES
    (
        '12345',
        'Ali Hassan',
        'ali@iba.edu.pk',
        '$2a$10$92IXUNpkjO0rOQ5byMi.Ye4oKoEa3Ro9llC/.og/at2.uheWG/igi',
        'student'
    ),
    (
        '22222',
        'Sara Khan',
        'sara@iba.edu.pk',
        '$2a$10$92IXUNpkjO0rOQ5byMi.Ye4oKoEa3Ro9llC/.og/at2.uheWG/igi',
        'student'
    );


-- Buildings
INSERT INTO
    buildings (id, name, location)
VALUES
    (
        '11111111-1111-1111-1111-111111111111',
        'Main Academic Block',
        'North Campus'
    ),
    (
        '22222222-2222-2222-2222-222222222222',
        'Business School',
        'East Wing'
    ),
    (
        '33333333-3333-3333-3333-333333333333',
        'Technology Block',
        'West Campus'
    );


-- Rooms
INSERT INTO
    rooms (building_id, name, capacity, type)
VALUES
    (
        '11111111-1111-1111-1111-111111111111',
        'Room 101',
        40,
        'Classroom'
    ),
    (
        '11111111-1111-1111-1111-111111111111',
        'Room 102',
        35,
        'Classroom'
    ),
    (
        '11111111-1111-1111-1111-111111111111',
        'Seminar Hall A',
        80,
        'Seminar Hall'
    ),
    (
        '11111111-1111-1111-1111-111111111111',
        'Computer Lab 1',
        30,
        'Computer Lab'
    ),
    (
        '22222222-2222-2222-2222-222222222222',
        'BS Conference Room',
        20,
        'Meeting Room'
    ),
    (
        '22222222-2222-2222-2222-222222222222',
        'Lecture Hall B1',
        60,
        'Classroom'
    ),
    (
        '33333333-3333-3333-3333-333333333333',
        'Tech Lab 01',
        25,
        'Computer Lab'
    ),
    (
        '33333333-3333-3333-3333-333333333333',
        'Innovation Hub',
        15,
        'Meeting Room'
    );


-- ── ROW LEVEL SECURITY (RLS) — disable for service role ──────
-- We use the service role key in the backend so RLS won't block us.
-- But enable it on all tables for safety:
ALTER TABLE users enable ROW level security;


ALTER TABLE buildings enable ROW level security;


ALTER TABLE rooms enable ROW level security;


ALTER TABLE bookings enable ROW level security;


ALTER TABLE blocked_slots enable ROW level security;


ALTER TABLE time_slots enable ROW level security;


-- Allow service role to bypass RLS (this is the default in Supabase)
-- No additional policies needed since we use service_role key server-side.

\end{lstlisting}

\subsection{seedTestData.sql (Data Injection)}
This script truncates existing data and injects the specific records required for the test cases defined in the Test Plan.
\begin{lstlisting}[language=SQL]
-- 1. DELETE all existing data (Order matters due to Foreign Keys)
-- Delete from child tables first, then parent tables
DELETE FROM bookings;


DELETE FROM blocked_slots;


DELETE FROM rooms;


DELETE FROM buildings;


DELETE FROM users;


-- 2. INSERT exactly these users
-- Password hash for 'testpass' using bcrypt (standard cost 10)
-- $2a$10$4r8D5gJGiqNuaPTR6qrxHeYlkUKl2nba97oZH.55EMivPEW8HoLe2
INSERT INTO
    users (erp, name, email, password, role)
VALUES
    (
        'test-student',
        'Test Student',
        'student@test.com',
        '$2a$10$4r8D5gJGiqNuaPTR6qrxHeYlkUKl2nba97oZH.55EMivPEW8HoLe2',
        'student'
    ),
    (
        'test-po',
        'Test PO',
        'po@test.com',
        '$2a$10$4r8D5gJGiqNuaPTR6qrxHeYlkUKl2nba97oZH.55EMivPEW8HoLe2',
        'programoffice'
    ),
    (
        'test-admin',
        'Test Admin',
        'admin@test.com',
        '$2a$10$4r8D5gJGiqNuaPTR6qrxHeYlkUKl2nba97oZH.55EMivPEW8HoLe2',
        'admin'
    ),
    (
        'test-student-2',
        'Concurrency Student',
        'student2@test.com',
        '$2a$10$4r8D5gJGiqNuaPTR6qrxHeYlkUKl2nba97oZH.55EMivPEW8HoLe2',
        'student'
    );


-- 3. INSERT one building
INSERT INTO
    buildings (id, name, location)
VALUES
    (
        '99999999-9999-4999-8999-999999999999',
        'Test Building',
        'Main Campus'
    );


-- 4. INSERT two rooms
INSERT INTO
    rooms (id, building_id, name, capacity, type)
VALUES
    (
        'aaaaaaaa-aaaa-4aaa-8aaa-aaaaaaaaaaaa',
        '99999999-9999-4999-8999-999999999999',
        'Test Room A',
        30,
        'Classroom'
    ),
    (
        'bbbbbbbb-bbbb-4bbb-8bbb-bbbbbbbbbbbb',
        '99999999-9999-4999-8999-999999999999',
        'Test Room B',
        20,
        'Meeting Room'
    );


-- 5. INSERT one existing booking for Test Room A, slot 1, a future date
-- We use current_date + interval '7 days' to ensure it's always in the future
INSERT INTO
    bookings (user_id, room_id, slot_id, date, purpose, status)
SELECT
    id,
    'aaaaaaaa-aaaa-4aaa-8aaa-aaaaaaaaaaaa',
    1,
    current_date + INTERVAL '7 days',
    'Existing Booking for Conflict Test',
    'approved'
FROM
    users
WHERE
    erp = 'test-student'
LIMIT
    1;
\end{lstlisting}

\subsection{seed.helper.ts (TypeScript Reset Helper)}
This function is called by the test runner to execute the reset logic programmatically.
\begin{lstlisting}[language=Java]
import { Client } from "pg";
import * as fs from "fs";
import * as path from "path";
import * as dotenv from "dotenv";

// Load the environment variables from .env.test
dotenv.config({ path: ".env.test" });

export async function resetDatabase() {
    // Use DATABASE_URL from .env.test
    const connectionString = process.env.DATABASE_URL;

    if (!connectionString) {
        console.error("Error: DATABASE_URL not found in .env.test");
        process.exit(1);
    }

    const client = new Client({
        connectionString,
        ssl: { rejectUnauthorized: false }, // Required for Supabase remote connections
    });

    try {
        await client.connect();
        console.log(`Connected to Supabase.`);

        const filePath = path.join(__dirname, "./../../test-data", "seedTestData.sql");

        console.log(`Reading file: ${filePath}`);
        const sql = fs.readFileSync(filePath, "utf8");

        console.log(`Executing SQL...`);
        await client.query(sql);

        console.log(`✅ Successfully executed seedTestData.sql`);
    } catch (err: any) {
        console.error(`❌ Error during seeding:`);
        console.error(err.message);
        process.exit(1);
    } finally {
        await client.end();
    }
}
\end{lstlisting}

\end{document}